%------------------------------------------------------------------------------%

\usepackage{apresentacao}
\usepackage{mathconfig}
\usepackage{tabto}

\makeatletter
\graphicspath  {{./figs/}}
\def\input@path{{./figs/}}
\makeatother

%------------------------------------------------------------------------------%
\title {Funções}
\author{\large Luis Alberto D'Afonseca}
\date  {Integrais e Séries}

%------------------------------------------------------------------------------%
\begin{document}

%------------------------------------------------------------------------------%
\begin{frame}
    \titlepage
\end{frame}

%------------------------------------------------------------------------------%
\section{Funções}

\begin{frame}{Função}
  \emph{Função} de $D$ em $C$
  \[
    \begin{array}{rccl}
      f\colon & D & \to & C \\
              & x & \to & y=f(x)
    \end{array}
  \]
  \pause
  associa a cada elemento, $x$, do \emph{domínio}, $D$, \\
  um único elemento, $y$, do \emph{contradomínio} $C$

  \vspace{\baselineskip}
  \pause
  \emph{Imagem} \\
  \pause
  Conjunto dos elementos $y$ associados a algum $x$
\end{frame}

%------------------------------------------------------------------------------%
\begin{frame}{Função Real}
  \emph{Função} de $\R$ em $\R$
  \[
  \begin{array}{rccl}
    f\colon & \R & \to & \R \\
    & x & \to & y=f(x)
  \end{array}
  \]
  domínio e contra domínio são $\R$
\end{frame}

%------------------------------------------------------------------------------%
\begin{frame}{Funções Comuns}
  \begin{description}[<+->]
    \item[Função constante]  \tabto{30mm} \(f(x) = c\)      \\[5mm]
    \item[Função afim]       \tabto{30mm} \(f(x) = ax + b\) \\[5mm]
    \item[Função potência]   \tabto{30mm} \(f(x) = x^n\)    \\[5mm]
    \item[Função polinômial] \tabto{30mm} \(f(x) = a_nx^n + a_{n-1}x^{n-1}
                                                 + \cdots + a_1 x + a_0\)
  \end{description}
\end{frame}

%------------------------------------------------------------------------------%
\begin{frame}{Funções Exponencial e Logarítmica}
  \begin{description}[<+->]
    \item[Função exponencial]          \tabto{30mm} \(f(x) = a^x\)       \\[5mm]
    \item[Função exponencial natural]  \tabto{30mm} \(f(x) = e^x\)       \\[5mm]
    \item[Função logarítmo]            \tabto{30mm} \(f(x) = \log_b(x)\) \\[5mm]
    \item[Função logarítmo natural]    \tabto{30mm} \(f(x) = \ln(x) = \log_e(x)\)
  \end{description}
  \onslide<+->{
  \begin{align*}
    y & = \log_b(x) & \quad\Rightarrow\qquad b^y = x & \hspace{90mm}~\\[3mm]
    y & = \ln(x)    & \quad\Rightarrow\qquad e^y = x &
  \end{align*}
  \vspace{-\baselineskip}
  }
\end{frame}

%------------------------------------------------------------------------------%
\begin{frame}{Funções Trigonométricas}
  \begin{description}[<+->]
    \item[Seno]       \tabto{30mm} \(f(x) = \sin(x)\)       \\[5mm]
    \item[Cosseno]    \tabto{30mm} \(f(x) = \cos(x)\)       \\[5mm]
    \item[Tangente]   \tabto{30mm} \(f(x) = \tan(x) = \frac{\sin(x)}{\cos(x)}\) \\[5mm]
    \item[Cotangente] \tabto{30mm} \(f(x) = \cot(x) = \frac{\cos(x)}{\sin(x)} = \frac{1}{\tan(x)}\) \\[5mm]
    \item[Secante]    \tabto{30mm} \(f(x) = \sec(x) = \frac{1}{\cos(x)} \)      \\[5mm]
    \item[Cossecante] \tabto{30mm} \(f(x) = \csc(x) = \frac{1}{\sin(x)} \)
  \end{description}
\end{frame}

%------------------------------------------------------------------------------%
\begin{frame}
  \tableofcontents
\end{frame}

%------------------------------------------------------------------------------%
\end{document}
%------------------------------------------------------------------------------%
